%% ============================================================================
%%  parked_sections.tex -- parked 2026-09-29. NOT part of the paper.
%%  Nothing \input's this file, and the whole body sits inside \iffalse ... \fi,
%%  so it stays inert even if someone \input's it by mistake.
%%
%%  What: text of the old draft (papers/main.tex, git 1029dd7, IEEEtran journal
%%  class) on active view selection (EFE) and on localization, copied verbatim.
%%  The full old draft is in git (git show 1029dd7:papers/main.tex).
%%
%%  Why (fact sheet 2026-09-29, section C,
%%  docs/reports/2026-09-29-paper-revision-1/logs/fact_sheet.md):
%%   * Localization: "Do NOT claim localization (the same group's concurrent
%%     localization paper covers it); remove the localization table, the ATE
%%     columns and the keyword."
%%   * Active view selection (EFE): "Remove the active-view-selection (EFE)
%%     section (no experiment); park removed text in papers/parked_sections.tex."
%%  The parked text also uses wording the fact sheet now forbids ("principled",
%%  "heuristic stable-object selection", "calibrated uncertainty", pixel mass
%%  n^eff, the r-tilde shrinkage). Rewrite it before any reuse.
%% ============================================================================
\iffalse

%% ----------------------------------------------------------------------------
%% 1. Related Work paragraph: visual localization (old main.tex lines 138-148)
%% ----------------------------------------------------------------------------
\textbf{Visual localization and rendering-based refinement.} Absolute pose
regression~\cite{chen2024marepo} and scene coordinate
regression~\cite{brachmann2023ace} learn a mapping from images to pose but
generalize poorly across scenes and degrade in complex layouts. Rendering-based
methods refine a coarse pose against a learned scene representation; replacing
NeRF~\cite{mildenhall2020nerf} with 3D-GS~\cite{kerbl2023} improves render
quality and efficiency for relocalization. \cite{chu2026semantic} augment a
Gaussian map with semantics and a language model to select stable objects for
robust indoor localization. We share the goal of robustness to scene change but
replace heuristic stable-object selection with a calibrated uncertainty derived
directly from the map.

%% ----------------------------------------------------------------------------
%% 2. Related Work paragraph: active perception (old main.tex lines 181-189)
%% ----------------------------------------------------------------------------
\textbf{Active perception.} Active inference casts perception and action as
minimizing expected free energy, which decomposes into information-seeking
(epistemic) and goal-directed (pragmatic) value~\cite{friston2010,
parr2022active}. Recent semantic 3D-GS exploration combines geometric and
semantic information gain with hand-set weights~\cite{nagami2025vista}. Our
disentangled uncertainty
provides the epistemic term in a principled form and, by gating out aleatoric
regions, avoids expending observations where more sensing cannot help
(Sec.~\ref{sec:efe}).

%% ----------------------------------------------------------------------------
%% 3. Method subsection: active viewpoint selection (old main.tex lines 371-386)
%% ----------------------------------------------------------------------------
\subsection{Active Viewpoint Selection (Optional)}
\label{sec:efe}
The disentangled uncertainty also yields a principled next-best-view objective.
Casting selection as minimizing expected free energy~\cite{friston2010}, a
candidate view $v$ with predicted per-Gaussian visibility $\nu_{v,n}$ is scored
by an epistemic term that seeks \emph{reducible} uncertainty and a pragmatic term
that seeks task-relevant content:
\begin{equation}
\mathrm{EFE}(v) = -\Big[\underbrace{\sum_n \nu_{v,n}\,u^{\mathrm{epi}}_n
(1-u^{\mathrm{ale}}_n)}_{\text{epistemic value}}
+ \lambda_p \underbrace{\sum_n \nu_{v,n}\,s_n}_{\text{pragmatic value}}\Big],
\end{equation}
where $s_n=\max(0,\muv^{\star\top}_n \mathbf{t})$ is similarity to a query
embedding $\mathbf{t}$. Gating the epistemic term by $(1-u^{\mathrm{ale}}_n)$
prevents the agent from repeatedly visiting genuinely ambiguous or dynamic
regions, a failure mode of objectives that do not separate the two uncertainties.

%% ----------------------------------------------------------------------------
%% 4. Experiments: localization setup sentences (old main.tex lines 400-413)
%% ----------------------------------------------------------------------------
%% From "Baselines":
For localization we compare against
representative APR~\cite{chen2024marepo}, SCR~\cite{brachmann2023ace}, and
GS-based relocalization methods, including the semantic localizer
of~\cite{chu2026semantic}. \todo{finalize baseline list and citations.}
%% From "Metrics":
localization by median translation/rotation error
(ATE/ARE) and success rate at $5\,$cm/$5^\circ$ and $2\,$cm/$2^\circ$;

%% ----------------------------------------------------------------------------
%% 5. Experiments subsection and table: localization under scene change
%%    (old main.tex lines 456-480)
%% ----------------------------------------------------------------------------
\subsection{Localization Under Scene Change}
Table~\ref{tab:loc} reports localization accuracy. The consolidation mechanism
(Sec.~\ref{sec:consolidation}) supplies a robust subset for matching under
rearrangement. \todo{interpret: e.g.\ that aleatoric-driven consolidation
preserves accuracy as the scene changes, versus baselines that degrade.}

\begin{table}[t]
\centering
\caption{Localization on dynamic indoor data (median ATE/ARE and success rate).
\todo{replace placeholders.}}
\label{tab:loc}
\begin{tabular}{lcccc}
\toprule
Method & ATE\,[cm]\,$\downarrow$ & ARE\,[\textdegree]\,$\downarrow$
& \multicolumn{2}{c}{Success\,[\%]\,$\uparrow$}\\
\cmidrule(lr){4-5}
 & & & 5cm/5\textdegree & 2cm/2\textdegree\\
\midrule
APR~\cite{chen2024marepo}        & \ph & \ph & \ph & \ph\\
SCR~\cite{brachmann2023ace}      & \ph & \ph & \ph & \ph\\
Semantic GS~\cite{chu2026semantic} & \ph & \ph & \ph & \ph\\
\textbf{Ours}                    & \textbf{\ph} & \textbf{\ph} & \textbf{\ph} & \textbf{\ph}\\
\bottomrule
\end{tabular}
\end{table}

%% ----------------------------------------------------------------------------
%% 6. Old ablation table: EFE row and ATE column (old main.tex lines 482-505)
%% ----------------------------------------------------------------------------
Table~\ref{tab:abl} isolates each component. \todo{interpret the four ablations:
(i) directional vMF vs.\ categorical posterior on open-vocabulary queries;
(ii) the 2-D $(n^{\mathrm{eff}},\bar r)$ uncertainty with the small-sample
correction vs.\ a single scalar; (iii) consolidation on/off under scene change;
(iv) the EFE objective vs.\ a hand-weighted geometric+semantic score.}

\begin{table}[t]
\centering
\caption{Ablation study. \todo{replace placeholders.}}
\label{tab:abl}
\begin{tabular}{lccc}
\toprule
Configuration & mIoU\,$\uparrow$ & AUSE$_{\mathrm{ale}}\!\downarrow$
& ATE\,[cm]\,$\downarrow$\\
\midrule
Categorical posterior        & \ph & \ph & \ph\\
Scalar uncertainty           & \ph & \ph & \ph\\
w/o consolidation            & \ph & \ph & \ph\\
Heuristic view score         & \ph & \ph & \ph\\
\textbf{Full model}          & \textbf{\ph} & \textbf{\ph} & \textbf{\ph}\\
\bottomrule
\end{tabular}
\end{table}

%% ----------------------------------------------------------------------------
%% 7. One-line mentions removed elsewhere (verbatim fragments)
%% ----------------------------------------------------------------------------
%% Abstract: "We evaluate calibration, open-vocabulary segmentation, and
%%   localization under scene change on indoor benchmarks."
%% Keywords: "Mapping, semantic scene understanding, localization, RGB-D
%%   perception, uncertainty estimation."
%% Introduction: "Such maps are attractive for localization and navigation
%%   because they combine appearance, geometry, and semantics in a single
%%   differentiable representation."
%% Introduction: "We evaluate calibration, open-vocabulary segmentation, and
%%   localization under scene change on standard indoor datasets, and ablate
%%   each component."
%% Pipeline caption: "... and used for consolidation (Sec.~\ref{sec:consolidation})
%%   and, optionally, active viewpoint selection (Sec.~\ref{sec:efe})."
%% Consolidation: "Conversely, Gaussians with $\bar r_n \geq r_{\mathrm{stable}}$
%%   are consolidated and define a \emph{robust} subset used for localization."

%% ----------------------------------------------------------------------------
%% 7. Parked 2026-09-29 in the fix pass (form review: page budget; claims
%%    review: the one-pass trick is prior work, joseph2025gradbp).
%%    Source: papers/main.tex before the fix pass.
%% ----------------------------------------------------------------------------
%% (a) Former fourth contribution bullet (demoted to an implementation note in
%%     Sec. III-C of main.tex):
\item \textbf{One backward pass for all statistics.} We follow the gradient
identity of~\cite{joseph2025gradbp} and stack feature, mass and depth channels
in one render, so that one backward pass gives every statistic. This needs a
rasterizer that is affine in the per-Gaussian features.
\todo{anonymous code link, or remove}

%% (b) Former Table IV (run time and memory), replaced by one paragraph
%%     ("Run time and memory") at the end of the Ablations subsection.
\subsection{Run Time and Memory}
\label{sec:timing}
\begin{table}[t]
\centering
\caption{Time per keyframe and peak GPU memory. One validation scene per
dataset, no other job running. \todo{hardware. Extractor name in row 1.
Replace placeholders.}}
\label{tab:timing}
\footnotesize
\setlength{\tabcolsep}{4pt}
\begin{tabular}{@{}lccc@{}}
\toprule
& \multicolumn{2}{c}{ms per keyframe$\downarrow$} & Memory\\
\cmidrule(lr){2-3}
Step & Replica & RIO10 & [GB]$\downarrow$\\
\midrule
Feature extraction                     & \ph & \ph & \ph\\
Render and backward pass               & \ph & \ph & \ph\\
Statistics, change test, re-admission  & \ph & \ph & \ph\\
\textbf{Ours, total}                   & \ph & \ph & \ph\\
\midrule
Averaging and NormLift, added          & \ph & \ph & \ph\\
CSS, added                             & \ph & \ph & \ph\\
\bottomrule
\end{tabular}
\end{table}


\fi
